\clearpage
\section*{Concluzii}
\addcontentsline{toc}{section}{Concluzii}

Scopul acestui proiect a fost analiza și modelarea unui sistem de gestiune a
librării online. Pe parcursul realizării acestui proiect am însușit limbajul
UML și am modelat o librărie online utilizând cunoștințele obținute. Am
realizat mai multe diagrame care expun diverse aspecte ale sistemului
respectiv.

Cu ajutorul limbajului UML putem modela o aplicație software pentru a ști exact
ce funcțional trebuie implmenetat și în ce mod. Am aplicat în practică
cunoștințele obținute și am realizat modelul aplicației BookMall.

Soft-ul folosit în cadrul acestui proiect a fost PlantUML. PlantUML permite
utilizatroului să creeze diagrame UML de toate tipurile și alte tipuri de
diagrame. Un comparație cu alte produse soft avantajul PlantUML este că
diagramele sunt create într-un fișier text dupa care este generat o imagine
pentru diagrama respectivă ceea ce este avantajos pentru persoanele care
preferă să nu folosească mousele și au un text editor preferat cum ar fi
Neovim.

De asemenea am început implementarea sistemului dat la moment fiind realizat
aproximativ 50\% din funcționalul stabilit de modelul UML. Am observat că în
unele cazuri nu este posibil de prevăzut din timp ce câmpuri va avea o clasă și
este nevoie apoi de actualizat modelul UML în conformitate cu cerințele reale.
Plănuiesc să finisez acest sistem fiind că este o tematică interesantă și până
la moment nu am mai dezvoltat astfel de sistem complex.

